\subsection{Plafonds empiriques du GH200 et dépendances locales}\label{sec:roofline-gpu}
Le modèle Roofline relie un débit de calcul à l'intensité du travail et à la bande passante disponible~\cite{roofline}. Pour NvPD, une cellule intérieure sans correspondance immédiate compare le coût d'édition au coût d'une correspondance antérieure. Le banc définit une mise à jour entière par l'expression \(\min(1+\min(u,d),\,v+(j-k))\), soit deux minima, deux additions et une soustraction sur des entiers de 32 bits. Le compilateur peut regrouper ces opérations dans une instruction machine.

Sur le GH200, une copie coalescée de 128\,Mio lus et 128\,Mio écrits donne \(B=3{,}488\pm0{,}007\)\,To/s. Un second noyau exécute huit chaînes indépendantes de mises à jour en registres, soit \(P=(5{,}481\pm0{,}013)\times10^{12}\) mises à jour par seconde. Les dispersions sont les écarts types de cinq processus. Pour \(I\) mises à jour par octet demandé, l'enveloppe empirique \(\min(BI,P)\) change de régime vers \(I=1{,}57\). Le banc de vérification lit trois mots et en écrit un par élément, puis varie le nombre de mises à jour réalisées en registres. Ses tableaux occupent 512\,Mio. Il atteint 214 milliards de mises à jour par seconde à \(I=1/16\), et 5\,068 milliards à \(I=16\).

Pour une colonne \(j\), la dernière occurrence \(k_s(j)\) d'un symbole \(s\) est croissante au sens large après sa première apparition. Sur un bloc commençant en \(L>0\), la plus petite colonne nécessaire sur la ligne précédente est donc bornée inférieurement par
\[
  \ell(L)=\min\bigl(\{L-1\}\cup\{k_s(L)-1\;:\;k_s(L)>0\}\bigr),\qquad \ell(0)=0.
\]
Si un symbole est absent jusqu'à \(L\), sa première occurrence ultérieure a un prédécesseur au moins égal à \(L\). En appliquant sept fois cette borne au bord de chaque bloc, le noyau peut produire huit lignes avec un seul lancement. Deux tableaux en mémoire partagée contiennent les lignes intermédiaires. Le noyau recalcule les colonnes du halo et écrit les cellules intérieures sur la dernière ligne. Le programme utilise la version par ligne lorsqu'un halo dépasse la capacité de mémoire partagée disponible.

\chart{gpu/roofline_gh200.png}{Le panneau a) utilise les octets demandés par le banc de calcul entier. Les panneaux b) et c) rapportent les cellules logiques \(mn\) au temps total de l'appel sur deux fenêtres homologues du chromosome~21. Chaque point repose sur cinq processus chauffés.}{fig:rooflinegpu}{bench/gpu/roofline/measure.py}

Les deux premières variantes produisent une ou deux lignes par lancement. La fusion recalcule jusqu'à trois valeurs intermédiaires. Une variante supplémentaire les partage entre les 32 threads d'un warp. Sur l'insertion de 131\,072 bases, la trace Nsight Systems compte 65\,537 lancements dans les deux cas. Le temps cumulé des noyaux passe de 154,87 à 162,97\,ms avec ces échanges. La fusion avec recalcul sert de témoin à la comparaison avec les tuiles locales.

Sur cette même fenêtre, les tuiles locales lancent 16\,385 noyaux, initialisation comprise, contre 131\,073 pour le traitement par ligne. Le temps complet descend de \(0{,}333\pm0{,}007\) à \(0{,}154\pm0{,}001\)\,s, soit un gain apparié de 2,16 (IC à 95\,\% de 2,12 à 2,20). Les huit entrées montrent un gain de 1,16 à 2,75 par rapport au traitement par ligne. Sur la fenêtre colinéaire de 524\,288 bases, le halo prend \(1{,}210\pm0{,}073\)\,s. La fenêtre avec insertion demande \(1{,}748\pm0{,}014\)\,s avec le halo et \(1{,}399\pm0{,}019\)\,s avec la fusion de deux lignes. Le plan de cette fenêtre présente un halo moyen de 77,4 colonnes et un maximum de 1\,804. Le double tampon est dimensionné pour le bloc le plus large.

Une dernière expérience répartit les blocs en deux groupes selon la largeur du halo. Sur la grande insertion, la mémoire partagée demandée descend de 22\,624 à 8\,960 octets pour 474 des 525 blocs. À nombre de cellules identique, les lancements passent de 65\,537 à 131\,073 et leur durée cumulée de 1,686 à 2,374\,s. Sur cette allocation, les temps complets sont de \(1{,}751\pm0{,}020\)\,s pour un seul groupe et de \(2{,}491\pm0{,}022\)\,s pour deux. Le rapport apparié vaut 0,703 (IC à 95\,\% de 0,697 à 0,709). Le regroupement ralentit les huit cas. Le mode public \texttt{nvpd-cuda-local} conserve un lancement par paquet de huit lignes.

Nsight Compute renvoie \texttt{ERR\_NVGPUCTRPERM} sur cette allocation. Le pilote réserve la lecture des compteurs matériels aux utilisateurs autorisés.\footnote{Voir le \href{https://developer.nvidia.com/ERR_NVGPUCTRPERM}{diagnostic de NVIDIA} et la \reposource{results/roofline/gpu/839191/ncu-local.txt}.} Les mesures d'activité permettent de vérifier les lancements et les durées des noyaux. Pour la grande insertion, les branches du code natif référencent entre 12 et 20 octets par cellule intérieure, soit 0,0557 cellule par octet demandé en moyenne sur le rectangle. Le trafic physique entre caches et HBM dépend du réemploi de ces adresses. La figure~\ref{fig:rooflinegpu} sépare les plafonds calibrés des débits complets de NvPD.

\synthese{Les tuiles locales réduisent la durée du calcul dense sur les huit fenêtres testées. Pour l'insertion de 524\,288 bases, la fusion de deux lignes reste la variante CUDA la plus rapide à 1,399\,s. Edlib et WFA2 calculent le même score de 13\,706 en 0,238 et 0,221\,s sur un cœur Grace.}
